\scanpage{26}

5. Modules libres de rang $2$ sur $\Lambda=\mathbb Z/4\mathbb Z$.

Soit $M$ un tel module, et soit
\[
M_0=M/2M=M\otimes_\Lambda\mathbb F_2 \tag*{(1)}
\]
qui est un vectoriel de rang $2$ sur $\mathbb F_2$. Soit $J=M_0^*$, on utilise la notation indicielle pour les éléments de $M_0^*$, notés $x_i$ $(i\in J)$. \deleted{Pour tout $i\in J$, soit}

On a une suite exacte de \deleted{groupes comm.} \inserted{$\Lambda$-modules}.
\diagram{\begin{tikzcd}[ampersand replacement=\&,column sep=small,row sep=small]
(2)\&0\arrow[r,no head]\&M_0\arrow[r,no head]\&M\arrow[r,no head,"\mathrm{can}"]\&M_0\arrow[r]\&0\\
\&\&M\arrow[u,dotted,"\mathrm{can}"]\arrow[ur,dotted,"2"]\&\&\&
\end{tikzcd}}
donc $M$ apparaît comme extension de $M_0$ par $M_0$ (en tant que $\Lambda$-module). Pour tt $x\in M_0$, soit $M_x$ son image inverse dans $M$, qui est un torseur sous $M_0$ (i.e. un ens. à quatre éléments \inserted{$M_x$}, muni d'un iso. $M_0\simeq V_{M_x}$).

On posera, pour $i\in J$
\[
\left\{\begin{aligned}
M_i&=M_{x_i}=\deleted{\unsure{\text{image inverse de }\{x_i\}}}\ \inserted{\text{fibre de }M\text{ sur }}x_i\in M_0,\ M_0\text{-torseur}\\
M_J&=\prod_{i\in J}M_i,\quad\text{un torseur sous }\deleted{\unsure{\mathbb F^J|\mathbb F}}\ M_0^J.
\end{aligned}\right. \tag*{(3)}
\]
Comme on a
\[
\sum_{i\in J}x_i=0 \tag*{(4)}
\]
la structure de groupe sur $M$ donne un iso. can.
\[
\bigwedge_{i\in J}M_i=0 \tag*{(6)}
\]
en d'autres termes, introduisant
\[
V_J(M_0)=\{(\deleted{\unsure{x}}\unsure{\xi|\zeta}_i)_{i\in J}\in M_0^J\mid\sum\unsure{\zeta|\xi}_i=0\} \tag*{(7)}
\]
s'insérant dans la suite exacte
\diagram{\begin{tikzcd}[ampersand replacement=\&,column sep=small]
(8)\&0\arrow[r,no head]\&V_J(M_0)\arrow[r,no head]\&M_0^J\arrow[r,no head,"\text{somme}"']\&M_0\arrow[r,no head]\&0,
\end{tikzcd}}
le torseur $M_J$ sous $M_0^J$ provient canoniquement d'un torseur sous $V_J(M_0)$, que je note $P_M$,

% archive sheet number: 25

